\documentclass[UTF8]{ctexart}
\usepackage{xeCJK}
\usepackage{geometry}
\geometry{left=2cm,right=2cm,top=2.5cm,bottom=2.5cm}
\setlength{\headheight}{12.7pt}

\usepackage{titlesec}
\titleformat{\section}
  {\normalfont\large\bfseries}
  {\thesection}
  {1em}
  {}
\titlespacing*{\section}
  {0pt}
  {3.5ex plus 1ex minus .2ex}
  {2.3ex plus .2ex}

\usepackage{amsmath}
\usepackage{amssymb}
\usepackage{amsthm}
\usepackage{enumitem}
\setlist[enumerate]{itemsep=0pt,topsep=0pt,parsep=0pt,leftmargin=0pt,itemindent=2.5em}
\setlist[itemize]{itemsep=0pt,topsep=0pt,parsep=0pt,leftmargin=0pt,itemindent=2.5em}
\usepackage{fancyhdr}
\pagestyle{fancy}
\fancyhf{}
\usepackage{graphicx}
\usepackage{hyperref}
\usepackage{indentfirst}
\usepackage{lmodern}


\begin{document}
\setlength{\abovedisplayskip}{1pt}
\setlength{\belowdisplayskip}{1pt}

\section{自然数的加法}

% 数学真正的起点.

\hypertarget{sec:定义1.1}{}
\noindent\textbf{定义1.1 (自然数的加法)}

给定\href{../01 自然数/01 自然数#nameddest=sec:定义1.3}{自然数} $m$,
递归定义映射$\mathrm{add}_m:\omega\to\omega$如下:
\begin{enumerate}
    \item $\mathrm{add}_m(0)\overset{\mathrm{def}}{=}m$,
    \item $\mathrm{add}_m(n^+)\overset{\mathrm{def}}{=}\mathrm{add}_m(n)^+$.
\end{enumerate}

接下来, 对任意的自然数$m,n$, 定义
\[
    m+n\overset{\mathrm{def}}{=}\mathrm{add}_m(n),\\[-3pt]
\]
称$+$为\textbf{加法}, 称$m+n$为$m,n$的\textbf{和}.

\vspace{10pt}

依定义, 任给自然数$n$, 有$n+1=\mathrm{add}_n(0^+)=\mathrm{add}_n(0)^+=n^+$,
即自然数的后继等同于加$1$.

% Soul Power !
% 2026.04.12

\vspace{10pt}

\hypertarget{sec:定理1.1}{}
\noindent\textbf{定理1.1 (自然数半群)}

自然数集上的加法有以下性质, 从而
是\href{../../../04 抽象代数/01 群论/01 群的基本概念/01 半群与群#nameddest=sec:定义1.1}
{交换幺半群}.
\begin{enumerate}
    \item 结合律: $\forall m,n,k\in\omega:\,(m+n)+k=m+(n+k)$,
    \item 单位元: $\forall n\in\omega:\,n+0=0+n=n$,
    \item 交换律: $\forall m,n\in\omega:\,m+n=n+m$.
\end{enumerate}

\noindent{\kaishu 证明.}

\begin{enumerate}
    \item 任给自然数$m,n$, 对$k$归纳证明$(m+n)+k=m+(n+k)$.
    
    \hspace{2.17em}
    首先$(m+n)+0=m+n=m+(n+0)$;
    
    \hspace{2.17em}
    其次任取自然数$k$, 假设$(m+n)+k=m+(n+k)$, 那么
    \[
        (m+n)+k^+=((m+n)+k)^+=(m+(n+k))^+=m+(n+k)^+=m+(n+k^+).
    \]

    \item 显然$\forall n\in\omega:\,n+0=n$. 再对$n$归纳证明$0+n=n$.
    
    \hspace{2.17em}
    首先$0+0=0$; 其次任取自然数$n$, 假设$0+n=n$, 那么
    \[
        0+n^+=(0+n)^+=n^+.
    \]

    \item 第一, 对$m$归纳证明$m+1=1+m$.
    
    \hspace{2.17em}
    首先$0+1=1+0$; 其次任取自然数$m$, 假设$m+1=1+m$, 那么
    \[
        m^++1=(m+1)+1=(1+m)+1=1+(m+1)=1+m^+.
    \]

    \hspace{2.17em}
    然后任给自然数$m$, 对$n$归纳证明$m+n=n+m$.

    \hspace{2.17em}
    首先$m+0=0+m$; 其次任取自然数$n$, 假设$m+n=n+m$, 那么
    \begin{equation*}
        m+n^+=(m+n)^+=(n+m)^+=n+(m+1)=n+(1+m)=(n+1)+m=n^++m.
    \tag*{\qedsymbol}\end{equation*}
\end{enumerate}





\newpage

自然数的加法和
它的\href{../01 自然数/03 自然数的序#nameddest=sec:定理1.4}{序}有严格相容性.

\vspace{10pt}

\hypertarget{sec:定理1.2}{}
\noindent\textbf{定理1.2}

任给自然数$m,n,k$, 如果$m<n$, 那么$m+k<n+k$.

\noindent{\kaishu 证明.}

给定自然数$m,n$, 假设$m<n$, 对$k$归纳证明$m+k<n+k$.

首先$m+0<n+0$ ; 其次任取自然数$k$, 假设$m+k<n+k$,
那么对任意的$x<(m+k)^+$有
\[
    x<(m+k)^+\implies x\leqslant m+k\implies
    x\leqslant n+k\implies x<(n+k)^+,
\]
从而$(m+k)^+\leqslant(n+k)^+$.

此时再假设$(m+k)^+=(n+k)^+$则有$m+k=n+k$矛盾, 所以$(m+k)^+<(n+k)^+$.
\hfill $\qedsymbol$

\vspace{10pt}

最后, 可以使用序关系来定义自然数的减法.

\vspace{10pt}

\hypertarget{sec:定义1.2}{}
\noindent\textbf{定义1.2 (自然数的减法)}

给定自然数$m,n$, 如果$m\leqslant n$, 那么存在唯一的自然数$q$使得
\[
    m+q=n,
\]
记作$n-m$, 称作$n,m$的\textbf{差}.

\noindent{\kaishu 验证.}

\noindent{\kaishu 存在性.}

给定自然数$m$, 对$n$归纳证明$n\geqslant m\to\exists q\in\omega:m+q=n$.

首先如果$0\geqslant m$, 那么$m=0$, 有$0+0=0$;

其次任取自然数$n$, 假设$n\geqslant m\to\exists q\in\omega:m+q=n$,
再假设$n^+\geqslant m$.

\begin{enumerate}
    \item 如果$m<n^+$, 那么$m\leqslant n$, 依假设存在自然数$q$使得
    $m+q=n$, 从而$m+q^+=n^+$.
    \item 如果$m=n^+$, 那么$m+0=n^+$.
\end{enumerate}
综上, 存在自然数$q'$使得$m+q'=n^+$.

\noindent{\kaishu 唯一性.}

假设同时存在自然数$q,q'$使得$m+q=m+q'=n$.

此时如果$q<q'$则$m+q<m+q'$, 矛盾. 同理$q'<q$也矛盾, 所以$q=q'$. \hfill $\qedsymbol$

\end{document}